\section{Revisão da Literatura}

% (i) teoria da produção e representação do fator ambiental
% (Cobb--Douglas, Solow e Dasgupta);

% (ii) externalidades, direitos de propriedade e internalização
% (Pigou e Coase);

% (iii) recursos de uso comum e governança institucional
% (Hardin e Ostrom);

% (iv) capital natural, sustentabilidade e incentivos intertemporais
% (Solow, Hartwick e Dasgupta);

% (v) síntese teórica e identificação da lacuna analítica.

% fim %